\begin{abstrak}{Maksimum 5 frasa kata, dipisahkan oleh tanda koma.}
Abstrak adalah versi ringkas dari tesis yang dituliskan dengan sistematis.
Abstrak seharusnya memberikan gambaran ringkas dari konten-konten yang penting
dari tesis. Kata-kata yang ada pada abstrak seharusnya dipilih dengan cermat
dan disusun dengan jelas untuk mendeskripsikan isi tesis. Abstrak memuat
pembahasan berikut: Latar belakang permasalahan (sekitar 1 kalimat),
Permasalahan (1-2 kalimat), Solusi yang ditawarkan (sekitar 5 kalimat), dan
Kinerja (sekitar 2 kalimat). Latar belakang permasalahan diambil dari problem
domain pada penelitian tesis. Permasalahan menjelaskan tingkat urgensi dari
problem pada penelitian tesis. Solusi mendiskusikan tujuan utama pada tesis,
klaim orisinalitas penelitian, skala penelitian, dan mendiskusikan metodologi
penyelesaian masalah. Kinerja membahas ringkasan eksperimen, hasil ujicoba dan
analisis kinerja.
\end{abstrak}

% Abstrak berbahasa Inggris (opsional):
% \begin{abstract}{keyword 1, keyword 2}
% ...
% \end{abstract}
